\documentclass[12pt,a4paper]{article}
% Окремий документ для перегляду результатів перевірки (XeLaTeX).
% Складання з каталогу research_nutrition:  xelatex glpk_appendix.tex  (двічі).
% Таблиці й абзаци можна також вставити в основний текст командами \input.
% Поля, кегль і пакети збігаються з diet_corrected.tex. Шрифт Liberation Serif має
% ті самі метрики, що й Nimbus Roman основного тексту; за потреби його можна замінити.
\usepackage{fontspec}
\setmainfont{Liberation Serif}
\newfontfamily\cyrillicfont{Liberation Serif}
\setmonofont{DejaVu Sans Mono}[Scale=MatchLowercase]
\newfontfamily\cyrillicfonttt{DejaVu Sans Mono}[Scale=MatchLowercase]
\usepackage{polyglossia}
\setdefaultlanguage{ukrainian}
\usepackage{amsmath,amssymb,booktabs,longtable,array,geometry}
\usepackage[section]{placeins}
\geometry{left=2.7cm,right=2cm,top=2cm,bottom=2.1cm}
\emergencystretch=4em
\setlength{\parskip}{3pt}
\renewcommand{\tablename}{Таблиця}
\begin{document}
\begin{center}
{\large\bfseries Перевірка обчислень задачі Стіглера в GNU Octave через GLPK}
\end{center}

% Згенеровано run_all.m. Абзаци з результатами перевірки в GNU Octave через GLPK.
% Початковий текст роботи не змінено; ці абзаци подано окремо.

\paragraph{Середовище та параметри розв’язувача.}
Обчислення виконано в GNU Octave 8.4.0 функцією \texttt{glpk}, яка викликає бібліотеку GLPK 5.0 (Ubuntu 24.04.5 LTS). Усі задачі розв’язано двофазним прямим симплекс-методом без пресолвера (\texttt{lpsolver}=1, \texttt{dual}=1, \texttt{presol}=0). Використано рівноважне масштабування (\texttt{scale}=16), вибір змінної за найкрутішим ребром (\texttt{price}=34) і тест відношень Гарріса (\texttt{rtest}=34). Допуски розв’язувача: \texttt{tolbnd}$=10^{-9}$, \texttt{toldj}$=10^{-9}$, \texttt{tolpiv}$=10^{-10}$. Запуск вважали успішним лише за коду завершення~0 і статусу~5 (оптимальний розв’язок). Цю умову виконано в усіх 294 запусках. Носій визначали умовою $x_i^*>10^{-9}$, активні рядки --- умовою $|s_j|\le10^{-9}\max\{1,|b_j|\}$.

\paragraph{Вихідні дані.}
Матрицю 1939 року взято з файлу \texttt{stigler.mod} архіву GLPK~4.65 без змін: контрольна сума SHA-256 файлу (\texttt{b0930f90e9b9ddee\ldots}) збігається із сумою оригіналу. Файл містить 77 продуктів і дев’ять показників у порядку: енергія, білок, кальцій, залізо, вітамін~A, тіамін, рибофлавін, ніацин, вітамін~C. Коефіцієнти задано на один долар витрат, ціни --- в центах за торгову одиницю. Для 58 продуктів із ваговою торговою одиницею перевірено стовпець їстівної маси на один долар: у 57 випадках він узгоджується з ціною та масою торгової одиниці. Для продукту \texttt{cornsirup} файл містить 4996~г, тоді як торгова одиниця 24 унції за 13{,}7~цента дає 4966{,}4~г; цей стовпець до моделі не входить. 

Матрицю 1944 року утворюють 14 рядків таблиці~B на с.~307 статті Стіглера. Таблиця~B не містить стовпця торгових одиниць: для продуктів №~1--74 одиниці взято з таблиці~A, для №~78--80 --- з приміток до таблиці~B. Порожні клітинки передано нулями. Звірка файлу \texttt{stigler1944.csv} зі сканом сторінки підтвердила 139 зі 140 клітинок (ціна та дев’ять показників у кожному рядку). Для 11 продуктів, які є в таблиці~A, перевірено співвідношення $a^{1944}_{ji}=a^{1939}_{ji}p^{1939}_i/p^{1944}_i$. Відхилення не перевищує 0{,}53 одиниці останнього друкованого знака в усіх клітинках, крім однієї. Виняток --- білок продукту №~3: у файлі 398~г, за таблицею~A --- 393{,}25~г; остання цифра на скані читається неоднозначно (3 або 8). Заміна 398 на 393 не змінює $C^*$, $x^*$, $u^*$, зведені вартості та значення $m_j$: цей продукт не входить до носія, а множник білкового рядка дорівнює нулю. Для продуктів №~78 і №~80 їстівна маса на один долар із приміток (4647 і 2071~г) узгоджується з ціною та торговою одиницею (4647{,}5 і 2071{,}2~г).

\paragraph{Пряма і двоїста задачі.}
Для кожної матриці окремо розв’язано пряму задачу $\min\{\mathbf1^\top x:\ Ax\ge b,\ x\ge0\}$ і двоїсту задачу $\max\{b^\top u:\ A^\top u\le\mathbf1,\ u\ge0\}$. Функція \texttt{glpk} повертає множники рядків як похідні оптимального значення за правими частинами. Тому для прямої задачі з рядками «$\ge$» множники збігаються з $u^*\ge0$, а для двоїстої задачі з рядками «$\le$» --- з $x^*\ge0$; зміна знака не потрібна. Множники прямої задачі відрізняються від розв’язку двоїстої не більше ніж на $7{,}63\cdot10^{-17}$ (1939) і $1{,}73\cdot10^{-17}$ (1944).

Одержано $C^*_{1939}=0{,}108662278206757$ і $C^*_{1944}=0{,}164110826270237$ дол./день. За 365 днів це 39{,}661732 і 59{,}900452 дол., за $365{,}25$ дня --- 39{,}688897 і 59{,}941479 дол. У табл.~\ref{tab:glpk-solution-1939} і \ref{tab:glpk-solution-1944} наведено оптимальні вектори $x^*$, у табл.~\ref{tab:glpk-dual-1939} і \ref{tab:glpk-dual-1944} --- надходження $Ax^*$, запаси $s=Ax^*-b$ і множники $u^*$. Повні вектори з 77 і 14 компонент та зведені вартості $r=\mathbf1-A^\top u^*$ збережено у файлах CSV. Порушення умов $Ax^*\ge b$, $x^*\ge0$, $A^\top u^*\le\mathbf1$, $u^*\ge0$ не перевищує $2{,}84\cdot10^{-14}$, двоїстий розрив --- $1{,}11\cdot10^{-16}$, добутки $|u_j^*s_j|$ і $|x_i^*r_i|$ --- $1{,}08\cdot10^{-16}$ (табл.~\ref{tab:glpk-checks}). В обох задачах додатними є п’ять компонент $x^*$ і п’ять компонент $u^*$. Активні рядки: енергія, кальцій, вітамін~A, вітамін~B$_2$, вітамін~C.

\paragraph{Фізичні кількості.}
Кількість торгових одиниць обчислено за формулою $z_i=100x_i^*/p_i$, де $p_i$ --- ціна торгової одиниці в центах. У табл.~\ref{tab:glpk-units-1939} і \ref{tab:glpk-units-1944} наведено добові кількості та річні кількості за 365 і $365{,}25$ дня. Торгова одиниця пшеничного борошна дорівнює 10 фунтам, борошна для млинців --- 20 унціям, тобто $1{,}25$ фунта. За $365{,}25$ дня розв’язок 1944 року містить 534{,}151 фунта пшеничного борошна і 134{,}414 фунта борошна для млинців, за 365 днів --- 533{,}786 і 134{,}322 фунта.

\paragraph{Історичний підсумок 1944 року.}
Сума вартостей у стовпці 1944 року таблиці~2 дорівнює 59{,}88 дол. Мінімум для опублікованої матриці становить 59{,}9005 дол. за 365 днів і 59{,}9415 дол. за $365{,}25$ дня. Стаття Стіглера не називає кількість днів у році; сума $24{,}50$ дол. на с.~311 відповідає 365 дням ($365\cdot3/44{,}7=24{,}50$). Тому історичний підсумок зіставлено зі значенням за 365 днів: різниця дорівнює 0{,}0205 дол., або 0{,}034\%. Історичний раціон містить ті самі п’ять продуктів, що й носій $x^*$ (табл.~\ref{tab:glpk-historic-1944}). Якщо його річні вартості поділити на 365, то за коефіцієнтами таблиці~B добове надходження становить 99{,}83\% вимоги для кальцію, 99{,}91\% вимоги для вітаміну~A, 99{,}82\% вимоги для вітаміну~B$_2$ (табл.~\ref{tab:glpk-historic-1944-nutrients}). Отже, історичний раціон не є допустимим для надрукованої матриці, і його вартість не суперечить мінімуму 59{,}9005 дол. За властивістю $V(\theta b)=\theta V(b)$ бюджет 59{,}88 дол. покриває щонайбільше 99{,}966\% усіх вимог одночасно. Коефіцієнти таблиці~B надруковано з кроком $0{,}1$ або 1. Якщо всі ненульові коефіцієнти збільшити на половину кроку, мінімальна вартість за 365 днів дорівнює 59{,}6902 дол.; якщо зменшити --- 60{,}1122 дол. Мінімум для кожної матриці, сумісної з округленням, лежить між цими значеннями, оскільки поелементне збільшення $A$ не збільшує $V$. Історичний підсумок 59{,}88 дол. належить цьому проміжку. Перевірка не відновлює неокруглені дані Стіглера.

\paragraph{Єдиність оптимального вектора.}
На множині $\{x:\ Ax\ge b,\ x\ge0,\ \mathbf1^\top x=C^*\}$ із повним обчисленим значенням $C^*$ знайдено мінімум і максимум кожної координати: 154 задачі для 1939 року і 28 задач для 1944 року. Усі задачі завершено з кодом~0 і статусом~5. Найбільша різниця між верхньою та нижньою межами координати дорівнює $5{,}95\cdot10^{-15}$ і $3{,}75\cdot10^{-16}$ дол./день (табл.~\ref{tab:glpk-uniqueness}). Це числовий результат, одержаний із допусками \texttt{tolbnd}$=$\texttt{toldj}$=10^{-9}$; він не є доведенням єдиності.

Строгий висновок одержано точною перевіркою сертифіката. Усі коефіцієнти мають щонайбільше один десятковий знак, тому $10A$ і $10b$ --- цілочислові. Для $S=\{i:x_i^*>0\}$ і $J_+=\{j:u_j^*>0\}$ у цілочисловій арифметиці без округлень обчислено визначник і приєднану матрицю для $10A_{J_+,S}$. Одержано $\det(10A^{1939}_{J_+,S})=-10457121500604640$ і $\det(10A^{1944}_{J_+,S})=5346316706566060$, тобто $\det A^{1939}_{J_+,S}=-65357009378779/625$ і $\det A^{1944}_{J_+,S}=267315835328303/5000$. Точно перевірено такі умови: $x_S^*>0$; $u_{J_+}^*>0$; $Ax^*\ge b$ з рівністю в рядках $J_+$; $A^\top u^*\le\mathbf1$ з рівністю в стовпцях $S$; $r_i>0$ для всіх $i\notin S$; $\mathbf1^\top x^*=b^\top u^*$. Найменша зведена вартість поза носієм дорівнює 0{,}04366642 для 1939 року і 0{,}09656077 для 1944 року. Отже, для обох табличних матриць оптимальний вектор $x^*$ єдиний. Запаси в рядках поза $J_+$ строго додатні, тому двоїстий оптимум $u^*$ також єдиний. Висновок стосується зазначених десяткових коефіцієнтів.

\paragraph{Чутливість до правих частин.}
Кожну з дев’яти вимог збільшено окремо на $1\%$ за незмінних інших. В усіх 18 випадках фактичний приріст вартості збігається з $0{,}01\,b_ju_j^*$; найбільша абсолютна різниця дорівнює $2{,}78\cdot10^{-17}$ дол. (табл.~\ref{tab:glpk-sensitivity}). Точна перевірка показала, що в кожному з цих випадків базис $(S,J_+)$ залишається оптимальним. Тому рівність $\Delta C^*=0{,}01\,b_ju_j^*$ для цих 18 збільшень є точною. Результат не поширюється на довільні зміни правих частин. За одночасного збільшення всіх вимог на $1\%$ одержано 0{,}109748900989 і 0{,}165751934533 дол./день. Приріст дорівнює $0{,}01\,C^*$ з різницею не більше ніж $1{,}11\cdot10^{-17}$ дол. і збігається із сумою дев’яти окремих приростів.

\paragraph{Неактивність і надлишковість.}
Для кожного неактивного рядка обчислено $m_j=\min\{A_{j\cdot}x:\ A_{k\cdot}x\ge b_k\ (k\ne j),\ x\ge0\}$ (табл.~\ref{tab:glpk-redundancy}). В усіх 8 випадках $m_j<b_j$, тому жодне неактивне обмеження не є надлишковим. Вилучення будь-якого одного рядка з $u_j^*=0$ і всіх чотирьох таких рядків одночасно змінює мінімальну вартість не більше ніж на $2{,}08\cdot10^{-16}$ дол./день, а оптимальний вектор --- не більше ніж на $7{,}89\cdot10^{-17}$ дол./день.

\paragraph{Семиденна задача.}
Сума семи окремих добових мінімумів дорівнює 0{,}796714750932 дол., мінімум спільної семиденної задачі --- 0{,}796714750932 дол.; різниця дорівнює $0$ (табл.~\ref{tab:glpk-multiday-days}). Добовий мінімум дорівнює 0{,}108662278207 дол. для вимог $(3;\ 70)$, 0{,}113403160673 дол. для $(3{,}3;\ 73{,}5)$ і 0{,}119590356250 дол. для $(3{,}6;\ 77)$. У дні 1 і 4 носій складається з п’яти продуктів. У дні 2, 3, 5, 6, 7 він складається з чотирьох продуктів (без яловичої печінки), а рядок вітаміну~B$_2$ неактивний. Кожний із семи добових оптимумів єдиний (точна перевірка сертифіката). Тому сукупність добових оптимумів єдина, а середні надходження в ній визначено однозначно.

Середнє надходження заліза за сім днів у цій сукупності дорівнює 62{,}039142467~мг. Порівняння в цілих числах показує, що воно строго більше за 12 і 62{,}03914~мг та строго менше за 65{,}14113~мг. За твердженням~1 межі 12 і 62{,}03914~мг зберігають $C^*$, а межа 65{,}14113~мг збільшує її. GLPK дає $C^*=0{,}796714751$ дол. для перших двох меж і $C^*=0{,}809187146$ дол. для третьої; приріст дорівнює 0{,}012472395 дол. (табл.~\ref{tab:glpk-multiday-iron}).

Для білка на множині $D_j=\{2,3,5,6,7\}$ середня добова нижня межа дорівнює 74{,}9~г, а середнє надходження --- 158{,}661312557~г. Воно строго більше за 74{,}9 і 158{,}66131~г та строго менше за 166{,}594~г. GLPK дає $C^*=0{,}796714751$ дол. для перших двох меж і $C^*=0{,}815475724$ дол. для третьої; приріст дорівнює 0{,}018760973 дол. (табл.~\ref{tab:glpk-multiday-protein}).

Допоміжні максимізації середнього надходження виконано за бюджету $0{,}79671475+10^{-7}$ дол. Найбільше середнє надходження заліза дорівнює 62{,}039167~мг, білка на множині $D_j$ --- 158{,}661354~г. Це результати з допуском до вартості: бюджет перевищує $C^*$ на $9{,}91\cdot10^{-8}$ дол. На точній оптимальній грані $\mathbf1^\top x=C^*$ мінімум і максимум середнього надходження збігаються: 62{,}039142~мг для заліза і 158{,}661313~г для білка.

\paragraph{Порівняння з текстом роботи.}
У табл.~\ref{tab:glpk-comparison} значення з тексту зіставлено з результатами GLPK; таблиця містить 164 рядки. Для 162 значень, які не залежать від розв’язувача, результат GLPK після округлення до наведеної в тексті кількості знаків збігається з текстом у 162 випадках. Числових розбіжностей не виявлено. Значення $1{,}23\cdot10^{-15}$ і $2{,}02\cdot10^{-16}$ у підрозд.~2.3 є числовими похибками іншого розв’язувача; GLPK дає $5{,}95\cdot10^{-15}$ і $3{,}75\cdot10^{-16}$.

\paragraph{Запропоновані виправлення.}
\begin{enumerate}
\item Підрозд.~1.1, «Обчислювальна перевірка». Речення про \texttt{glpsol} версії~4.65, SciPy і SymPy описує попередній запуск. Для результатів цього додатка середовище таке: GNU Octave 8.4.0, бібліотека GLPK 5.0; версія~4.65 стосується лише файлу даних \texttt{stigler.mod}.
\item Підрозд.~2.3. Якщо різниці меж координат наводяться для GLPK, значення $1{,}23\cdot10^{-15}$ і $2{,}02\cdot10^{-16}$ слід замінити на $5{,}95\cdot10^{-15}$ і $3{,}75\cdot10^{-16}$.
\item Файл \texttt{stigler1944.csv}, продукт №~3, білок. Значення 398 не підтверджено однозначно; за таблицею~A очікується 393. На числа розділу ця клітинка не впливає.
\item Підрозд.~1.1, зіставлення з підсумком $59{,}88$. Текст порівнює його зі значенням за 365 днів (59{,}9005 дол.); це узгоджується зі статтею. Доцільно додати обґрунтування: стаття не називає кількість днів, але сума $24{,}50$ дол. на с.~311 відповідає 365 дням.
\end{enumerate}


\clearpage
% Згенеровано run_all.m (GNU Octave, GLPK). Потребує пакета booktabs.
\begin{table}[htbp]
\centering\small
\caption{Оптимальний вектор витрат, 1939 рік (GLPK); решта компонент дорівнюють нулю}
\label{tab:glpk-solution-1939}
\begin{tabular}{lrrr}
\toprule
Продукт & $x_i^*$, дол./день & Дол. за 365 днів & Дол. за $365{,}25$ дня \\
\midrule
Пшеничне борошно & 0{,}029519061676 & 10{,}774458 & 10{,}781837 \\
Яловича печінка & 0{,}001892557291 & 0{,}690783 & 0{,}691257 \\
Капуста & 0{,}011214435246 & 4{,}093269 & 4{,}096072 \\
Шпинат & 0{,}005007660467 & 1{,}827796 & 1{,}829048 \\
Сушена біла квасоля & 0{,}061028563527 & 22{,}275426 & 22{,}290683 \\
\midrule
Разом & 0{,}108662278207 & 39{,}661732 & 39{,}688897 \\
\bottomrule
\end{tabular}
\end{table}

% Згенеровано run_all.m (GNU Octave, GLPK). Потребує пакета booktabs.
\begin{table}[htbp]
\centering\footnotesize\setlength{\tabcolsep}{3pt}
\caption{Фізичні кількості продуктів, 1939 рік: $z_i=100x_i^*/p_i$ торгових одиниць на день}
\label{tab:glpk-units-1939}
\begin{tabular}{llrrrrrr}
\toprule
 & Торгова & $p_i$, & $z_i$, & \multicolumn{2}{c}{Торгових одиниць за} & \multicolumn{2}{c}{Фунтів за} \\
\cmidrule(lr){5-6}\cmidrule(lr){7-8}
Продукт & одиниця & центів & од./день & 365 днів & $365{,}25$ дня & 365 днів & $365{,}25$ дня \\
\midrule
Пшеничне борошно & 10 фунтів & 36{,}0 & 0{,}081997394 & 29{,}929049 & 29{,}949548 & 299{,}290486 & 299{,}495480 \\
Яловича печінка & 1 фунт & 26{,}8 & 0{,}007061781 & 2{,}577550 & 2{,}579315 & 2{,}577550 & 2{,}579315 \\
Капуста & 1 фунт & 3{,}7 & 0{,}303092844 & 110{,}628888 & 110{,}704661 & 110{,}628888 & 110{,}704661 \\
Шпинат & 1 фунт & 8{,}1 & 0{,}061822969 & 22{,}565384 & 22{,}580839 & 22{,}565384 & 22{,}580839 \\
Сушена біла квасоля & 1 фунт & 5{,}9 & 1{,}034382433 & 377{,}549588 & 377{,}808184 & 377{,}549588 & 377{,}808184 \\
\bottomrule
\end{tabular}
\end{table}

% Згенеровано run_all.m (GNU Octave, GLPK). Потребує пакета booktabs.
\begin{table}[htbp]
\centering\small
\caption{Надходження, запаси і двоїсті змінні, 1939 рік (GLPK)}
\label{tab:glpk-dual-1939}
\begin{tabular}{lrrrr}
\toprule
Показник & $b_j$ & $(Ax^*)_j$ & $s_j$ & $u_j^*$ \\
\midrule
Енергія, тис. ккал & 3{,}0 & 3{,}000000 & 0{,}000000 & 0{,}008765147 \\
Білок, г & 70{,}0 & 147{,}413535 & 77{,}413535 & 0{,}000000000 \\
Кальцій, г & 0{,}8 & 0{,}800000 & 0{,}000000 & 0{,}031737713 \\
Залізо, мг & 12{,}0 & 60{,}466922 & 48{,}466922 & 0{,}000000000 \\
Вітамін A, тис. МО & 5{,}0 & 5{,}000000 & 0{,}000000 & 0{,}000400233 \\
Вітамін B$_1$, мг & 1{,}8 & 4{,}120439 & 2{,}320439 & 0{,}000000000 \\
Вітамін B$_2$, мг & 2{,}7 & 2{,}700000 & 0{,}000000 & 0{,}016358033 \\
Ніацин, мг & 18{,}0 & 27{,}315981 & 9{,}315981 & 0{,}000000000 \\
Вітамін C, мг & 75{,}0 & 75{,}000000 & 0{,}000000 & 0{,}000144118 \\
\bottomrule
\end{tabular}
\end{table}

% Згенеровано run_all.m (GNU Octave, GLPK). Потребує пакета booktabs.
\begin{table}[htbp]
\centering\small
\caption{Оптимальний вектор витрат, 1944 рік (GLPK); решта компонент дорівнюють нулю}
\label{tab:glpk-solution-1944}
\begin{tabular}{lrrr}
\toprule
Продукт & $x_i^*$, дол./день & Дол. за 365 днів & Дол. за $365{,}25$ дня \\
\midrule
Пшеничне борошно & 0{,}094472740525 & 34{,}482550 & 34{,}506168 \\
Капуста & 0{,}014315196735 & 5{,}225047 & 5{,}228626 \\
Шпинат & 0{,}004249219142 & 1{,}550965 & 1{,}552027 \\
Борошно для млинців & 0{,}035917299511 & 13{,}109814 & 13{,}118794 \\
Свиняча печінка & 0{,}015156370357 & 5{,}532075 & 5{,}535864 \\
\midrule
Разом & 0{,}164110826270 & 59{,}900452 & 59{,}941479 \\
\bottomrule
\end{tabular}
\end{table}

% Згенеровано run_all.m (GNU Octave, GLPK). Потребує пакета booktabs.
\begin{table}[htbp]
\centering\footnotesize\setlength{\tabcolsep}{3pt}
\caption{Фізичні кількості продуктів, 1944 рік: $z_i=100x_i^*/p_i$ торгових одиниць на день}
\label{tab:glpk-units-1944}
\begin{tabular}{llrrrrrr}
\toprule
 & Торгова & $p_i$, & $z_i$, & \multicolumn{2}{c}{Торгових одиниць за} & \multicolumn{2}{c}{Фунтів за} \\
\cmidrule(lr){5-6}\cmidrule(lr){7-8}
Продукт & одиниця & центів & од./день & 365 днів & $365{,}25$ дня & 365 днів & $365{,}25$ дня \\
\midrule
Пшеничне борошно & 10 фунтів & 64{,}6 & 0{,}146242632 & 53{,}378561 & 53{,}415121 & 533{,}785608 & 534{,}151215 \\
Капуста & 1 фунт & 4{,}9 & 0{,}292146872 & 106{,}633608 & 106{,}706645 & 106{,}633608 & 106{,}706645 \\
Шпинат & 1 фунт & 11{,}6 & 0{,}036631200 & 13{,}370388 & 13{,}379546 & 13{,}370388 & 13{,}379546 \\
Борошно для млинців & 20 унцій & 12{,}2 & 0{,}294404094 & 107{,}457494 & 107{,}531095 & 134{,}321868 & 134{,}413869 \\
Свиняча печінка & 1 фунт & 21{,}9 & 0{,}069207171 & 25{,}260617 & 25{,}277919 & 25{,}260617 & 25{,}277919 \\
\bottomrule
\end{tabular}
\end{table}

% Згенеровано run_all.m (GNU Octave, GLPK). Потребує пакета booktabs.
\begin{table}[htbp]
\centering\small
\caption{Надходження, запаси і двоїсті змінні, 1944 рік (GLPK)}
\label{tab:glpk-dual-1944}
\begin{tabular}{lrrrr}
\toprule
Показник & $b_j$ & $(Ax^*)_j$ & $s_j$ & $u_j^*$ \\
\midrule
Енергія, тис. ккал & 3{,}0 & 3{,}000000 & 0{,}000000 & 0{,}028876611 \\
Білок, г & 70{,}0 & 99{,}303830 & 29{,}303830 & 0{,}000000000 \\
Кальцій, г & 0{,}8 & 0{,}800000 & 0{,}000000 & 0{,}028312380 \\
Залізо, мг & 12{,}0 & 29{,}475597 & 17{,}475597 & 0{,}000000000 \\
Вітамін A, тис. МО & 5{,}0 & 5{,}000000 & 0{,}000000 & 0{,}000724485 \\
Вітамін B$_1$, мг & 1{,}8 & 3{,}324068 & 1{,}524068 & 0{,}000000000 \\
Вітамін B$_2$, мг & 2{,}7 & 2{,}700000 & 0{,}000000 & 0{,}013431654 \\
Ніацин, мг & 18{,}0 & 32{,}250746 & 14{,}250746 & 0{,}000000000 \\
Вітамін C, мг & 75{,}0 & 75{,}000000 & 0{,}000000 & 0{,}000199243 \\
\bottomrule
\end{tabular}
\end{table}

% Згенеровано run_all.m (GNU Octave, GLPK). Потребує пакета booktabs.
\begin{table}[htbp]
\centering\small
\caption{Перевірка допустимості, двоїстого розриву та доповняльної нежорсткості}
\label{tab:glpk-checks}
\begin{tabular}{lrrr}
\toprule
Величина & 1939 & 1944 & Допуск \\
\midrule
Код завершення, статус (пряма задача) & 0; 5 & 0; 5 & 0; 5 \\
Код завершення, статус (двоїста задача) & 0; 5 & 0; 5 & 0; 5 \\
Порушення $Ax^*\ge b$, $x^*\ge0$ & $2{,}84\cdot10^{-14}$ & $4{,}44\cdot10^{-16}$ & $1{,}00\cdot10^{-9}$ \\
Порушення $A^\top u^*\le\mathbf1$, $u^*\ge0$ & $0$ & $4{,}44\cdot10^{-16}$ & $1{,}00\cdot10^{-9}$ \\
$|\mathbf1^\top x^*-b^\top u^*|$ & $1{,}11\cdot10^{-16}$ & $0$ & $1{,}00\cdot10^{-12}$ \\
Різниця значень прямої та двоїстої задач & $5{,}55\cdot10^{-17}$ & $1{,}11\cdot10^{-16}$ & $1{,}00\cdot10^{-12}$ \\
$\max_j|u_j^*s_j|$ & $5{,}81\cdot10^{-17}$ & $1{,}28\cdot10^{-17}$ & $1{,}00\cdot10^{-12}$ \\
$\max_i|x_i^*r_i|$ & $1{,}08\cdot10^{-16}$ & $3{,}99\cdot10^{-18}$ & $1{,}00\cdot10^{-12}$ \\
$\max_j|u_j^*-u_j^{\mathrm{D}}|$ & $7{,}63\cdot10^{-17}$ & $1{,}73\cdot10^{-17}$ & $1{,}00\cdot10^{-9}$ \\
$\max_i|x_i^*-\lambda_i^{\mathrm{D}}|$ & $7{,}63\cdot10^{-17}$ & $9{,}71\cdot10^{-17}$ & $1{,}00\cdot10^{-9}$ \\
\bottomrule
\end{tabular}
\end{table}

% Згенеровано run_all.m (GNU Octave, GLPK). Потребує пакета booktabs.
\begin{table}[htbp]
\centering\small
\caption{Узгодженість таблиці B з таблицею A: найбільше відхилення $a^{1944}_{ji}-a^{1939}_{ji}p^{1939}_i/p^{1944}_i$ в одиницях останнього друкованого знака}
\label{tab:glpk-datacheck-1944}
\begin{tabular}{rlrrr}
\toprule
№ & Продукт & $p_i^{1939}$, центів & $p_i^{1944}$, центів & Найбільше відхилення \\
\midrule
1 & Wheat flour & 36{,}0 & 64{,}6 & 0{,}43 \\
3 & Wheat cereal & 24{,}2 & 23{,}2 & 4{,}75 \\
5 & Corn meal & 4{,}6 & 6{,}3 & 0{,}41 \\
8 & Rolled oats & 7{,}1 & 9{,}9 & 0{,}53 \\
15 & Evaporated milk & 6{,}7 & 10{,}0 & 0{,}45 \\
46 & Cabbage & 3{,}7 & 4{,}9 & 0{,}39 \\
51 & Potatoes & 34{,}0 & 80{,}1 & 0{,}49 \\
52 & Spinach & 8{,}1 & 11{,}6 & 0{,}36 \\
53 & Sweet potatoes & 5{,}1 & 12{,}3 & 0{,}41 \\
69 & Navy beans & 5{,}9 & 10{,}8 & 0{,}45 \\
74 & Sugar & 51{,}7 & 67{,}0 & 0{,}30 \\
78 & Pancake flour & --- & 12{,}2 & --- \\
79 & Beets & --- & 7{,}3 & --- \\
80 & Pork liver & --- & 21{,}9 & --- \\
\bottomrule
\end{tabular}
\end{table}

% Згенеровано run_all.m (GNU Octave, GLPK). Потребує пакета booktabs.
\begin{table}[htbp]
\centering\small
\caption{Історичний раціон 1944 року і розв’язок ЛП для опублікованої матриці (365 днів)}
\label{tab:glpk-historic-1944}
\begin{tabular}{lrrrr}
\toprule
 & \multicolumn{2}{c}{Кількість, фунтів} & \multicolumn{2}{c}{Вартість, дол.} \\
\cmidrule(lr){2-3}\cmidrule(lr){4-5}
Продукт & Стіглер & ЛП & Стіглер & ЛП \\
\midrule
Пшеничне борошно & 535 & 533{,}79 & 34{,}53 & 34{,}4826 \\
Капуста & 107 & 106{,}63 & 5{,}23 & 5{,}2250 \\
Шпинат & 13 & 13{,}37 & 1{,}56 & 1{,}5510 \\
Борошно для млинців & 134 & 134{,}32 & 13{,}08 & 13{,}1098 \\
Свиняча печінка & 25 & 25{,}26 & 5{,}48 & 5{,}5321 \\
\midrule
Разом & & & 59{,}88 & 59{,}9005 \\
\bottomrule
\end{tabular}
\end{table}

% Згенеровано run_all.m (GNU Octave, GLPK). Потребує пакета booktabs.
\begin{table}[htbp]
\centering\small
\caption{Добове надходження історичного раціону 1944 року за опублікованою матрицею (річні вартості поділено на 365)}
\label{tab:glpk-historic-1944-nutrients}
\begin{tabular}{lrrr}
\toprule
Показник & $b_j$ & Надходження & Частка від $b_j$, \% \\
\midrule
Енергія, тис. ккал & 3{,}0 & 3{,}0016 & 100{,}05 \\
Білок, г & 70{,}0 & 99{,}3118 & 141{,}87 \\
Кальцій, г & 0{,}8 & 0{,}7987 & 99{,}83 \\
Залізо, мг & 12{,}0 & 29{,}4271 & 245{,}23 \\
Вітамін A, тис. МО & 5{,}0 & 4{,}9953 & 99{,}91 \\
Вітамін B$_1$, мг & 1{,}8 & 3{,}3265 & 184{,}81 \\
Вітамін B$_2$, мг & 2{,}7 & 2{,}6952 & 99{,}82 \\
Ніацин, мг & 18{,}0 & 32{,}2129 & 178{,}96 \\
Вітамін C, мг & 75{,}0 & 75{,}0199 & 100{,}03 \\
\bottomrule
\end{tabular}
\end{table}

% Згенеровано run_all.m (GNU Octave, GLPK). Потребує пакета booktabs.
\begin{table}[htbp]
\centering\small
\caption{Перевірка єдиності оптимального вектора}
\label{tab:glpk-uniqueness}
\begin{tabular}{lrr}
\toprule
Величина & 1939 & 1944 \\
\midrule
Кількість допоміжних задач & 154 & 28 \\
Задач із кодом 0 і статусом 5 & 154 & 28 \\
$\max_i(\overline x_i-\underline x_i)$, дол./день & $5{,}95\cdot10^{-15}$ & $3{,}75\cdot10^{-16}$ \\
Найбільше відхилення меж від $x_i^*$, дол./день & $4{,}93\cdot10^{-15}$ & $3{,}63\cdot10^{-16}$ \\
$|S|=|J_+|$ & 5 & 5 \\
$\det(10A_{J_+,S})$ & $-10457121500604640$ & $5346316706566060$ \\
$\min_{i\notin S}r_i$ & 0{,}0436664195 & 0{,}0965607671 \\
$\min_{j\notin J_+}s_j$ & 2{,}320439 & 1{,}524068 \\
Точна перевірка умов сертифіката & так & так \\
\bottomrule
\end{tabular}
\end{table}

% Згенеровано run_all.m (GNU Octave, GLPK). Потребує пакета booktabs.
\begin{table}[htbp]
\centering\footnotesize
\caption{Приріст мінімальної добової вартості при збільшенні вимог на $1\%$ і прогноз $0{,}01\,b_ju_j^*$, дол.}
\label{tab:glpk-sensitivity}
\begin{tabular}{lrrrr}
\toprule
 & \multicolumn{2}{c}{1939} & \multicolumn{2}{c}{1944} \\
\cmidrule(lr){2-3}\cmidrule(lr){4-5}
Змінена вимога & $\Delta C^*$ & $0{,}01\,b_ju_j^*$ & $\Delta C^*$ & $0{,}01\,b_ju_j^*$ \\
\midrule
Енергія, тис. ккал & 0{,}000262954419 & 0{,}000262954419 & 0{,}000866298340 & 0{,}000866298340 \\
Білок, г & 0{,}000000000000 & 0{,}000000000000 & 0{,}000000000000 & 0{,}000000000000 \\
Кальцій, г & 0{,}000253901708 & 0{,}000253901708 & 0{,}000226499039 & 0{,}000226499039 \\
Залізо, мг & 0{,}000000000000 & 0{,}000000000000 & 0{,}000000000000 & 0{,}000000000000 \\
Вітамін A, тис. МО & 0{,}000020011636 & 0{,}000020011636 & 0{,}000036224227 & 0{,}000036224227 \\
Вітамін B$_1$, мг & 0{,}000000000000 & 0{,}000000000000 & 0{,}000000000000 & 0{,}000000000000 \\
Вітамін B$_2$, мг & 0{,}000441666883 & 0{,}000441666883 & 0{,}000362654652 & 0{,}000362654652 \\
Ніацин, мг & 0{,}000000000000 & 0{,}000000000000 & 0{,}000000000000 & 0{,}000000000000 \\
Вітамін C, мг & 0{,}000108088137 & 0{,}000108088137 & 0{,}000149432005 & 0{,}000149432005 \\
\midrule
Усі дев’ять вимог & 0{,}001086622782 & 0{,}001086622782 & 0{,}001641108263 & 0{,}001641108263 \\
\bottomrule
\end{tabular}
\end{table}

% Згенеровано run_all.m (GNU Octave, GLPK). Потребує пакета booktabs.
\begin{table}[htbp]
\centering\small
\caption{Мінімальне надходження $m_j$ за решти обмежень для неактивних рядків}
\label{tab:glpk-redundancy}
\begin{tabular}{lrrr}
\toprule
Показник & $b_j$ & $m_j$, 1939 & $m_j$, 1944 \\
\midrule
Білок, г & 70{,}0 & 1{,}505268 & 37{,}730293 \\
Залізо, мг & 12{,}0 & 1{,}898087 & 11{,}008698 \\
Вітамін B$_1$, мг & 1{,}8 & 0{,}039391 & 0{,}826300 \\
Ніацин, мг & 18{,}0 & 0{,}000000 & 3{,}812587 \\
\bottomrule
\end{tabular}
\end{table}

% Згенеровано run_all.m (GNU Octave, GLPK). Потребує пакета booktabs.
\begin{table}[htbp]
\centering\small
\caption{Одноденні мінімуми та добові витрати у спільній семиденній задачі}
\label{tab:glpk-multiday-days}
\begin{tabular}{crrrr}
\toprule
День & Енергія, тис. ккал & Білок, г & $C_t^*$, дол. & У семиденній задачі, дол. \\
\midrule
1 & 3{,}0 & 70{,}0 & 0{,}108662278207 & 0{,}108662278207 \\
2 & 3{,}3 & 73{,}5 & 0{,}113403160673 & 0{,}113403160673 \\
3 & 3{,}6 & 77{,}0 & 0{,}119590356250 & 0{,}119590356250 \\
4 & 3{,}0 & 70{,}0 & 0{,}108662278207 & 0{,}108662278207 \\
5 & 3{,}3 & 73{,}5 & 0{,}113403160673 & 0{,}113403160673 \\
6 & 3{,}6 & 77{,}0 & 0{,}119590356250 & 0{,}119590356250 \\
7 & 3{,}3 & 73{,}5 & 0{,}113403160673 & 0{,}113403160673 \\
\midrule
Разом & & & 0{,}796714750932 & 0{,}796714750932 \\
\bottomrule
\end{tabular}
\end{table}

% Згенеровано run_all.m (GNU Octave, GLPK). Потребує пакета booktabs.
\begin{table}[htbp]
\centering\small
\caption{Вплив нижньої межі середнього надходження заліза за сім днів}
\label{tab:glpk-multiday-iron}
\begin{tabular}{rrrr}
\toprule
$B_j$, мг & $C^*$, дол. & $\Delta C^*$, дол. & Середнє в оптимумі, мг \\
\midrule
12{,}00000 & 0{,}796714750932 & 0{,}000000000000 & 62{,}03914247 \\
62{,}03914 & 0{,}796714750932 & 0{,}000000000000 & 62{,}03914247 \\
65{,}14113 & 0{,}809187146386 & 0{,}012472395454 & 65{,}14113000 \\
\bottomrule
\end{tabular}
\end{table}

% Згенеровано run_all.m (GNU Octave, GLPK). Потребує пакета booktabs.
\begin{table}[htbp]
\centering\small
\caption{Вплив нижньої межі середнього надходження білка на множині $D_j=\{2,3,5,6,7\}$}
\label{tab:glpk-multiday-protein}
\begin{tabular}{rrrr}
\toprule
$B_j$, г & $C^*$, дол. & $\Delta C^*$, дол. & Середнє в оптимумі, г \\
\midrule
74{,}90000 & 0{,}796714750932 & 0{,}000000000000 & 158{,}66131256 \\
158{,}66131 & 0{,}796714750932 & 0{,}000000000000 & 158{,}66131256 \\
166{,}59400 & 0{,}815475723898 & 0{,}018760972966 & 166{,}59400000 \\
\bottomrule
\end{tabular}
\end{table}

\FloatBarrier
\clearpage
% Згенеровано run_all.m; потребує пакетів longtable, booktabs, array.
{\scriptsize\setlength{\tabcolsep}{3pt}
\begin{longtable}{>{\raggedright}p{4.3cm}rrr>{\raggedright\arraybackslash}p{2.3cm}}
\caption{Порівняння значень у тексті роботи з результатами GLPK}\label{tab:glpk-comparison}\\
\toprule
Місце в роботі & Значення в тексті & Результат GLPK & Різниця & Висновок \\
\midrule
\endfirsthead
\toprule
Місце в роботі & Значення в тексті & Результат GLPK & Різниця & Висновок \\
\midrule
\endhead
\bottomrule
\endfoot
Підрозд. 1.1, обчислювальна перевірка: $C^*$(1939), дол./день & $0{,}108662278206757$ & $0{,}108662278206757$ & $-2{,}08\cdot10^{-16}$ & підтверджено \\
Підрозд. 1.1, обчислювальна перевірка: 365{,}25 $C^*$(1939), дол. & $39{,}688897115018$ & $39{,}68889711501792$ & $-8{,}53\cdot10^{-14}$ & підтверджено \\
Підрозд. 1.1, обчислювальна перевірка: 365 $C^*$(1939), дол. & $39{,}661731545466$ & $39{,}66173154546623$ & $2{,}34\cdot10^{-13}$ & підтверджено \\
Табл. 2: сума вартостей рядків, 1939 & $39{,}93$ & $39{,}9300$ & $7{,}11\cdot10^{-15}$ & підтверджено \\
Табл. 2: сума вартостей рядків, 1944 & $59{,}88$ & $59{,}8800$ & $7{,}11\cdot10^{-15}$ & підтверджено \\
Табл. 3: пшеничне борошно, фунтів за 365{,}25 дня & $299{,}50$ & $299{,}4955$ & $-4{,}52\cdot10^{-3}$ & підтверджено \\
Табл. 3: яловича печінка, фунтів за 365{,}25 дня & $2{,}58$ & $2{,}5793$ & $-6{,}85\cdot10^{-4}$ & підтверджено \\
Табл. 3: капуста, фунтів за 365{,}25 дня & $110{,}70$ & $110{,}7047$ & $4{,}66\cdot10^{-3}$ & підтверджено \\
Табл. 3: шпинат, фунтів за 365{,}25 дня & $22{,}58$ & $22{,}5808$ & $8{,}39\cdot10^{-4}$ & підтверджено \\
Табл. 3: біла квасоля, фунтів за 365{,}25 дня & $377{,}81$ & $377{,}8082$ & $-1{,}82\cdot10^{-3}$ & підтверджено \\
Табл. 3: річна вартість ЛП, дол. & $39{,}69$ & $39{,}6889$ & $-1{,}10\cdot10^{-3}$ & підтверджено \\
Підрозд. 1.1, після табл. 3: різниця 39{,}93 - 365{,}25 $C^*$, дол. & $0{,}24$ & $0{,}2411$ & $1{,}10\cdot10^{-3}$ & підтверджено \\
Табл. 4: пшеничне борошно, $x^*$, дол./день & $0{,}029519061676$ & $0{,}02951906167649$ & $4{,}88\cdot10^{-13}$ & підтверджено \\
Табл. 4: пшеничне борошно, фунтів за 365{,}25 дня & $299{,}495480$ & $299{,}49547993$ & $-7{,}40\cdot10^{-8}$ & підтверджено \\
Табл. 4: яловича печінка, $x^*$, дол./день & $0{,}001892557291$ & $0{,}00189255729071$ & $-2{,}95\cdot10^{-13}$ & підтверджено \\
Табл. 4: яловича печінка, фунтів за 365{,}25 дня & $2{,}579315$ & $2{,}57931549$ & $4{,}87\cdot10^{-7}$ & підтверджено \\
Табл. 4: капуста, $x^*$, дол./день & $0{,}011214435246$ & $0{,}01121443524614$ & $1{,}45\cdot10^{-13}$ & підтверджено \\
Табл. 4: капуста, фунтів за 365{,}25 дня & $110{,}704661$ & $110{,}70466145$ & $4{,}50\cdot10^{-7}$ & підтверджено \\
Табл. 4: шпинат, $x^*$, дол./день & $0{,}005007660467$ & $0{,}00500766046673$ & $-2{,}75\cdot10^{-13}$ & підтверджено \\
Табл. 4: шпинат, фунтів за 365{,}25 дня & $22{,}580839$ & $22{,}58083933$ & $3{,}27\cdot10^{-7}$ & підтверджено \\
Табл. 4: біла квасоля, $x^*$, дол./день & $0{,}061028563527$ & $0{,}06102856352669$ & $-3{,}07\cdot10^{-13}$ & підтверджено \\
Табл. 4: біла квасоля, фунтів за 365{,}25 дня & $377{,}808184$ & $377{,}80818353$ & $-4{,}72\cdot10^{-7}$ & підтверджено \\
Табл. 5: Енергія, тис. ккал: $(Ax^*)_j$ & $3{,}000000$ & $3{,}00000000$ & $-4{,}44\cdot10^{-16}$ & підтверджено \\
Табл. 5: Енергія, тис. ккал: $s_j$ & $0{,}000000$ & $0{,}00000000$ & $-4{,}44\cdot10^{-16}$ & підтверджено \\
Табл. 5: Енергія, тис. ккал: $u_j^*$ & $0{,}008765147$ & $0{,}00876514730$ & $2{,}98\cdot10^{-10}$ & підтверджено \\
Табл. 5: Білок, г: $(Ax^*)_j$ & $147{,}413535$ & $147{,}41353494$ & $-5{,}78\cdot10^{-8}$ & підтверджено \\
Табл. 5: Білок, г: $s_j$ & $77{,}413535$ & $77{,}41353494$ & $-5{,}78\cdot10^{-8}$ & підтверджено \\
Табл. 5: Білок, г: $u_j^*$ & $0{,}000000000$ & $0{,}00000000000$ & $0$ & підтверджено \\
Табл. 5: Кальцій, г: $(Ax^*)_j$ & $0{,}800000$ & $0{,}80000000$ & $2{,}22\cdot10^{-16}$ & підтверджено \\
Табл. 5: Кальцій, г: $s_j$ & $0{,}000000$ & $0{,}00000000$ & $2{,}22\cdot10^{-16}$ & підтверджено \\
Табл. 5: Кальцій, г: $u_j^*$ & $0{,}031737713$ & $0{,}03173771345$ & $4{,}46\cdot10^{-10}$ & підтверджено \\
Табл. 5: Залізо, мг: $(Ax^*)_j$ & $60{,}466922$ & $60{,}46692210$ & $1{,}02\cdot10^{-7}$ & підтверджено \\
Табл. 5: Залізо, мг: $s_j$ & $48{,}466922$ & $48{,}46692210$ & $1{,}02\cdot10^{-7}$ & підтверджено \\
Табл. 5: Залізо, мг: $u_j^*$ & $0{,}000000000$ & $0{,}00000000000$ & $0$ & підтверджено \\
Табл. 5: Вітамін A, тис. МО: $(Ax^*)_j$ & $5{,}000000$ & $5{,}00000000$ & $8{,}88\cdot10^{-16}$ & підтверджено \\
Табл. 5: Вітамін A, тис. МО: $s_j$ & $0{,}000000$ & $0{,}00000000$ & $8{,}88\cdot10^{-16}$ & підтверджено \\
Табл. 5: Вітамін A, тис. МО: $u_j^*$ & $0{,}000400233$ & $0{,}00040023272$ & $-2{,}78\cdot10^{-10}$ & підтверджено \\
Табл. 5: Вітамін B1, мг: $(Ax^*)_j$ & $4{,}120439$ & $4{,}12043880$ & $-1{,}95\cdot10^{-7}$ & підтверджено \\
Табл. 5: Вітамін B1, мг: $s_j$ & $2{,}320439$ & $2{,}32043880$ & $-1{,}95\cdot10^{-7}$ & підтверджено \\
Табл. 5: Вітамін B1, мг: $u_j^*$ & $0{,}000000000$ & $0{,}00000000000$ & $0$ & підтверджено \\
Табл. 5: Вітамін B2, мг: $(Ax^*)_j$ & $2{,}700000$ & $2{,}70000000$ & $-3{,}55\cdot10^{-15}$ & підтверджено \\
Табл. 5: Вітамін B2, мг: $s_j$ & $0{,}000000$ & $0{,}00000000$ & $-3{,}55\cdot10^{-15}$ & підтверджено \\
Табл. 5: Вітамін B2, мг: $u_j^*$ & $0{,}016358033$ & $0{,}01635803270$ & $-3{,}01\cdot10^{-10}$ & підтверджено \\
Табл. 5: Ніацин, мг: $(Ax^*)_j$ & $27{,}315981$ & $27{,}31598070$ & $-3{,}00\cdot10^{-7}$ & підтверджено \\
Табл. 5: Ніацин, мг: $s_j$ & $9{,}315981$ & $9{,}31598070$ & $-3{,}00\cdot10^{-7}$ & підтверджено \\
Табл. 5: Ніацин, мг: $u_j^*$ & $0{,}000000000$ & $0{,}00000000000$ & $0$ & підтверджено \\
Табл. 5: Вітамін C, мг: $(Ax^*)_j$ & $75{,}000000$ & $75{,}00000000$ & $-2{,}84\cdot10^{-14}$ & підтверджено \\
Табл. 5: Вітамін C, мг: $s_j$ & $0{,}000000$ & $0{,}00000000$ & $-2{,}84\cdot10^{-14}$ & підтверджено \\
Табл. 5: Вітамін C, мг: $u_j^*$ & $0{,}000144118$ & $0{,}00014411752$ & $-4{,}85\cdot10^{-10}$ & підтверджено \\
Підрозд. 1.1, розрахунок 1944 року: $C^*$(1944), дол./день & $0{,}164110826270237$ & $0{,}164110826270237$ & $3{,}33\cdot10^{-16}$ & підтверджено \\
Підрозд. 1.1, розрахунок 1944 року: 365{,}25 $C^*$(1944), дол. & $59{,}941479295204$ & $59{,}94147929520418$ & $1{,}85\cdot10^{-13}$ & підтверджено \\
Підрозд. 1.1, розрахунок 1944 року: 365 $C^*$(1944), дол. & $59{,}900451588637$ & $59{,}90045158863662$ & $-3{,}77\cdot10^{-13}$ & підтверджено \\
Табл. 6: пшеничне борошно, $x^*$, дол./день & $0{,}094472740525$ & $0{,}09447274052526$ & $2{,}64\cdot10^{-13}$ & підтверджено \\
Табл. 6: пшеничне борошно, фунтів за 365{,}25 дня & $534{,}151215$ & $534{,}15121481$ & $-1{,}88\cdot10^{-7}$ & підтверджено \\
Табл. 6: капуста, $x^*$, дол./день & $0{,}014315196735$ & $0{,}01431519673493$ & $-6{,}74\cdot10^{-14}$ & підтверджено \\
Табл. 6: капуста, фунтів за 365{,}25 дня & $106{,}706645$ & $106{,}70664505$ & $4{,}97\cdot10^{-8}$ & підтверджено \\
Табл. 6: шпинат, $x^*$, дол./день & $0{,}004249219142$ & $0{,}00424921914214$ & $1{,}43\cdot10^{-13}$ & підтверджено \\
Табл. 6: шпинат, фунтів за 365{,}25 дня & $13{,}379546$ & $13{,}37954562$ & $-3{,}82\cdot10^{-7}$ & підтверджено \\
Табл. 6: борошно для млинців, $x^*$, дол./день & $0{,}035917299511$ & $0{,}03591729951083$ & $-1{,}71\cdot10^{-13}$ & підтверджено \\
Табл. 6: борошно для млинців, фунтів за 365{,}25 дня & $134{,}413869$ & $134{,}41386933$ & $3{,}27\cdot10^{-7}$ & підтверджено \\
Табл. 6: свиняча печінка, $x^*$, дол./день & $0{,}015156370357$ & $0{,}01515637035707$ & $6{,}89\cdot10^{-14}$ & підтверджено \\
Табл. 6: свиняча печінка, фунтів за 365{,}25 дня & $25{,}277919$ & $25{,}27791905$ & $5{,}44\cdot10^{-8}$ & підтверджено \\
Табл. 7: Енергія, тис. ккал: $(Ax^*)_j$ & $3{,}000000$ & $3{,}00000000$ & $4{,}44\cdot10^{-16}$ & підтверджено \\
Табл. 7: Енергія, тис. ккал: $s_j$ & $0{,}000000$ & $0{,}00000000$ & $4{,}44\cdot10^{-16}$ & підтверджено \\
Табл. 7: Енергія, тис. ккал: $u_j^*$ & $0{,}028876611$ & $0{,}02887661132$ & $3{,}19\cdot10^{-10}$ & підтверджено \\
Табл. 7: Білок, г: $(Ax^*)_j$ & $99{,}303830$ & $99{,}30383033$ & $3{,}34\cdot10^{-7}$ & підтверджено \\
Табл. 7: Білок, г: $s_j$ & $29{,}303830$ & $29{,}30383033$ & $3{,}34\cdot10^{-7}$ & підтверджено \\
Табл. 7: Білок, г: $u_j^*$ & $0{,}000000000$ & $0{,}00000000000$ & $0$ & підтверджено \\
Табл. 7: Кальцій, г: $(Ax^*)_j$ & $0{,}800000$ & $0{,}80000000$ & $0$ & підтверджено \\
Табл. 7: Кальцій, г: $s_j$ & $0{,}000000$ & $0{,}00000000$ & $0$ & підтверджено \\
Табл. 7: Кальцій, г: $u_j^*$ & $0{,}028312380$ & $0{,}02831237993$ & $-7{,}10\cdot10^{-11}$ & підтверджено \\
Табл. 7: Залізо, мг: $(Ax^*)_j$ & $29{,}475597$ & $29{,}47559730$ & $2{,}99\cdot10^{-7}$ & підтверджено \\
Табл. 7: Залізо, мг: $s_j$ & $17{,}475597$ & $17{,}47559730$ & $2{,}99\cdot10^{-7}$ & підтверджено \\
Табл. 7: Залізо, мг: $u_j^*$ & $0{,}000000000$ & $0{,}00000000000$ & $0$ & підтверджено \\
Табл. 7: Вітамін A, тис. МО: $(Ax^*)_j$ & $5{,}000000$ & $5{,}00000000$ & $1{,}78\cdot10^{-15}$ & підтверджено \\
Табл. 7: Вітамін A, тис. МО: $s_j$ & $0{,}000000$ & $0{,}00000000$ & $1{,}78\cdot10^{-15}$ & підтверджено \\
Табл. 7: Вітамін A, тис. МО: $u_j^*$ & $0{,}000724485$ & $0{,}00072448453$ & $-4{,}67\cdot10^{-10}$ & підтверджено \\
Табл. 7: Вітамін B1, мг: $(Ax^*)_j$ & $3{,}324068$ & $3{,}32406816$ & $1{,}57\cdot10^{-7}$ & підтверджено \\
Табл. 7: Вітамін B1, мг: $s_j$ & $1{,}524068$ & $1{,}52406816$ & $1{,}57\cdot10^{-7}$ & підтверджено \\
Табл. 7: Вітамін B1, мг: $u_j^*$ & $0{,}000000000$ & $0{,}00000000000$ & $0$ & підтверджено \\
Табл. 7: Вітамін B2, мг: $(Ax^*)_j$ & $2{,}700000$ & $2{,}70000000$ & $-4{,}44\cdot10^{-16}$ & підтверджено \\
Табл. 7: Вітамін B2, мг: $s_j$ & $0{,}000000$ & $0{,}00000000$ & $-4{,}44\cdot10^{-16}$ & підтверджено \\
Табл. 7: Вітамін B2, мг: $u_j^*$ & $0{,}013431654$ & $0{,}01343165378$ & $-2{,}25\cdot10^{-10}$ & підтверджено \\
Табл. 7: Ніацин, мг: $(Ax^*)_j$ & $32{,}250746$ & $32{,}25074623$ & $2{,}33\cdot10^{-7}$ & підтверджено \\
Табл. 7: Ніацин, мг: $s_j$ & $14{,}250746$ & $14{,}25074623$ & $2{,}33\cdot10^{-7}$ & підтверджено \\
Табл. 7: Ніацин, мг: $u_j^*$ & $0{,}000000000$ & $0{,}00000000000$ & $0$ & підтверджено \\
Табл. 7: Вітамін C, мг: $(Ax^*)_j$ & $75{,}000000$ & $75{,}00000000$ & $0$ & підтверджено \\
Табл. 7: Вітамін C, мг: $s_j$ & $0{,}000000$ & $0{,}00000000$ & $0$ & підтверджено \\
Табл. 7: Вітамін C, мг: $u_j^*$ & $0{,}000199243$ & $0{,}00019924267$ & $-3{,}27\cdot10^{-10}$ & підтверджено \\
Табл. 8: Білок, г: $m_j$, 1939 & $1{,}505268$ & $1{,}50526844$ & $4{,}40\cdot10^{-7}$ & підтверджено \\
Табл. 8: Білок, г: $m_j$, 1944 & $37{,}730293$ & $37{,}73029293$ & $-6{,}84\cdot10^{-8}$ & підтверджено \\
Табл. 8: Залізо, мг: $m_j$, 1939 & $1{,}898087$ & $1{,}89808734$ & $3{,}41\cdot10^{-7}$ & підтверджено \\
Табл. 8: Залізо, мг: $m_j$, 1944 & $11{,}008698$ & $11{,}00869757$ & $-4{,}30\cdot10^{-7}$ & підтверджено \\
Табл. 8: Вітамін B1, мг: $m_j$, 1939 & $0{,}039391$ & $0{,}03939076$ & $-2{,}44\cdot10^{-7}$ & підтверджено \\
Табл. 8: Вітамін B1, мг: $m_j$, 1944 & $0{,}826300$ & $0{,}82630011$ & $1{,}06\cdot10^{-7}$ & підтверджено \\
Табл. 8: Ніацин, мг: $m_j$, 1939 & $0{,}000000$ & $0{,}00000000$ & $0$ & підтверджено \\
Табл. 8: Ніацин, мг: $m_j$, 1944 & $3{,}812587$ & $3{,}81258694$ & $-6{,}20\cdot10^{-8}$ & підтверджено \\
Табл. 10: день 1: окрема одноденна задача, дол. & $0{,}109$ & $0{,}10866$ & $-3{,}38\cdot10^{-4}$ & підтверджено \\
Табл. 10: день 1: день у семиденній задачі, дол. & $0{,}109$ & $0{,}10866$ & $-3{,}38\cdot10^{-4}$ & підтверджено \\
Табл. 10: день 2: окрема одноденна задача, дол. & $0{,}113$ & $0{,}11340$ & $4{,}03\cdot10^{-4}$ & підтверджено \\
Табл. 10: день 2: день у семиденній задачі, дол. & $0{,}113$ & $0{,}11340$ & $4{,}03\cdot10^{-4}$ & підтверджено \\
Табл. 10: день 3: окрема одноденна задача, дол. & $0{,}120$ & $0{,}11959$ & $-4{,}10\cdot10^{-4}$ & підтверджено \\
Табл. 10: день 3: день у семиденній задачі, дол. & $0{,}120$ & $0{,}11959$ & $-4{,}10\cdot10^{-4}$ & підтверджено \\
Табл. 10: день 4: окрема одноденна задача, дол. & $0{,}109$ & $0{,}10866$ & $-3{,}38\cdot10^{-4}$ & підтверджено \\
Табл. 10: день 4: день у семиденній задачі, дол. & $0{,}109$ & $0{,}10866$ & $-3{,}38\cdot10^{-4}$ & підтверджено \\
Табл. 10: день 5: окрема одноденна задача, дол. & $0{,}113$ & $0{,}11340$ & $4{,}03\cdot10^{-4}$ & підтверджено \\
Табл. 10: день 5: день у семиденній задачі, дол. & $0{,}113$ & $0{,}11340$ & $4{,}03\cdot10^{-4}$ & підтверджено \\
Табл. 10: день 6: окрема одноденна задача, дол. & $0{,}120$ & $0{,}11959$ & $-4{,}10\cdot10^{-4}$ & підтверджено \\
Табл. 10: день 6: день у семиденній задачі, дол. & $0{,}120$ & $0{,}11959$ & $-4{,}10\cdot10^{-4}$ & підтверджено \\
Табл. 10: день 7: окрема одноденна задача, дол. & $0{,}113$ & $0{,}11340$ & $4{,}03\cdot10^{-4}$ & підтверджено \\
Табл. 10: день 7: день у семиденній задачі, дол. & $0{,}113$ & $0{,}11340$ & $4{,}03\cdot10^{-4}$ & підтверджено \\
Табл. 10: разом: сума добових мінімумів, дол. & $0{,}797$ & $0{,}79671$ & $-2{,}85\cdot10^{-4}$ & підтверджено \\
Табл. 10: разом: семиденна задача, дол. & $0{,}797$ & $0{,}79671$ & $-2{,}85\cdot10^{-4}$ & підтверджено \\
Підрозд. 1.2, після табл. 10: модуль різниці $C^*$ і суми $C_t^*$ менший за $10^{-12}$ & $<1{,}00\cdot10^{-12}$ & $0$ & --- & підтверджено (нерівність виконується) \\
Підрозд. 1.2, після табл. 10: середнє надходження заліза, мг & $62{,}03914$ & $62{,}0391425$ & $2{,}47\cdot10^{-6}$ & підтверджено \\
Підрозд. 1.2, після табл. 10: середнє добових меж заліза, мг & $12$ & $12{,}00$ & $0$ & підтверджено \\
Підрозд. 1.2, після табл. 10: найбільше середнє заліза за бюджету 0{,}79671475+$10^{-7}$, мг & $62{,}03917$ & $62{,}0391671$ & $-2{,}89\cdot10^{-6}$ & підтверджено \\
Підрозд. 1.2, після табл. 10: 1{,}05 $\times$ найбільше середнє заліза, мг & $65{,}14113$ & $65{,}1411255$ & $-4{,}54\cdot10^{-6}$ & підтверджено \\
Табл. 11: B = 12.00000 мг: $C^*$, дол. & $0{,}796715$ & $0{,}79671475$ & $-2{,}49\cdot10^{-7}$ & підтверджено \\
Табл. 11: B = 12.00000 мг: приріст $C^*$, дол. & $0{,}000000$ & $0{,}00000000$ & $1{,}11\cdot10^{-15}$ & підтверджено \\
Табл. 11: B = 62.03914 мг: $C^*$, дол. & $0{,}796715$ & $0{,}79671475$ & $-2{,}49\cdot10^{-7}$ & підтверджено \\
Табл. 11: B = 62.03914 мг: приріст $C^*$, дол. & $0{,}000000$ & $0{,}00000000$ & $6{,}66\cdot10^{-16}$ & підтверджено \\
Табл. 11: B = 65.14113 мг: $C^*$, дол. & $0{,}809187$ & $0{,}80918715$ & $1{,}46\cdot10^{-7}$ & підтверджено \\
Табл. 11: B = 65.14113 мг: приріст $C^*$, дол. & $0{,}012472$ & $0{,}01247240$ & $3{,}95\cdot10^{-7}$ & підтверджено \\
Підрозд. 1.3, обчислювальна перевірка: вартість семиденної задачі, дол. & $0{,}79671475$ & $0{,}7967147509$ & $9{,}32\cdot10^{-10}$ & підтверджено \\
Підрозд. 1.3, обчислювальна перевірка: середня добова межа білка на $D_j$, г & $74{,}9$ & $74{,}900$ & $0$ & підтверджено \\
Підрозд. 1.3, обчислювальна перевірка: середнє надходження білка на $D_j$, г & $158{,}66131$ & $158{,}6613126$ & $2{,}56\cdot10^{-6}$ & підтверджено \\
Підрозд. 1.3, обчислювальна перевірка: найбільше середнє білка за бюджету 0{,}79671475+$10^{-7}$, г & $158{,}66135$ & $158{,}6613544$ & $4{,}45\cdot10^{-6}$ & підтверджено \\
Підрозд. 1.3, обчислювальна перевірка: 1{,}05 $\times$ найбільше середнє білка, г & $166{,}594$ & $166{,}59442$ & $4{,}22\cdot10^{-4}$ & підтверджено \\
Табл. 12: межа 74.90000 г: вартість, дол. & $0{,}796715$ & $0{,}79671475$ & $-2{,}49\cdot10^{-7}$ & підтверджено \\
Табл. 12: межа 74.90000 г: приріст, дол. & $0{,}000000$ & $0{,}00000000$ & $2{,}78\cdot10^{-15}$ & підтверджено \\
Табл. 12: межа 158.66131 г: вартість, дол. & $0{,}796715$ & $0{,}79671475$ & $-2{,}49\cdot10^{-7}$ & підтверджено \\
Табл. 12: межа 158.66131 г: приріст, дол. & $0{,}000000$ & $0{,}00000000$ & $1{,}11\cdot10^{-16}$ & підтверджено \\
Табл. 12: межа 166.59400 г: вартість, дол. & $0{,}815476$ & $0{,}81547572$ & $-2{,}76\cdot10^{-7}$ & підтверджено \\
Табл. 12: межа 166.59400 г: приріст, дол. & $0{,}018761$ & $0{,}01876097$ & $-2{,}70\cdot10^{-8}$ & підтверджено \\
Підрозд. 1.3, після табл. 12: вартість за межі 166{,}594 г, дол. & $0{,}81547572$ & $0{,}8154757239$ & $3{,}90\cdot10^{-9}$ & підтверджено \\
Підрозд. 1.3, після табл. 12: приріст за межі 166{,}594 г, дол. & $0{,}01876097$ & $0{,}0187609730$ & $2{,}97\cdot10^{-9}$ & підтверджено \\
Підрозд. 2.3: найбільша різниця меж координат, 1939 (154 задачі) & $1{,}23\cdot10^{-15}$ & $5{,}95\cdot10^{-15}$ & $4{,}72\cdot10^{-15}$ & числова похибка, залежить від розв’язувача; для GLPK: $5{,}95\cdot10^{-15}$ \\
Підрозд. 2.3: найбільша різниця меж координат, 1944 (28 задач) & $2{,}02\cdot10^{-16}$ & $3{,}75\cdot10^{-16}$ & $1{,}73\cdot10^{-16}$ & числова похибка, залежить від розв’язувача; для GLPK: $3{,}75\cdot10^{-16}$ \\
Підрозд. 2.4: det $A_{J_+,S}$, 1939 & $-65357009378779/625$ & $-10457121500604640/100000$ & $0$ & підтверджено точно (цілочислова арифметика) \\
Підрозд. 2.4: det $A_{J_+,S}$, 1944 & $267315835328303/5000$ & $5346316706566060/100000$ & $0$ & підтверджено точно (цілочислова арифметика) \\
Підрозд. 2.4: найменша зведена вартість поза носієм, 1939 & $0{,}04366642$ & $0{,}0436664195$ & $-4{,}58\cdot10^{-10}$ & підтверджено \\
Підрозд. 2.4: найменша зведена вартість поза носієм, 1944 & $0{,}09656077$ & $0{,}0965607671$ & $-2{,}90\cdot10^{-9}$ & підтверджено \\
Табл. 13: Енергія, тис. ккал: приріст $C^*$, 1939 & $0{,}000262954$ & $0{,}00026295442$ & $4{,}19\cdot10^{-10}$ & підтверджено \\
Табл. 13: Енергія, тис. ккал: приріст $C^*$, 1944 & $0{,}000866298$ & $0{,}00086629834$ & $3{,}40\cdot10^{-10}$ & підтверджено \\
Табл. 13: Білок, г: приріст $C^*$, 1939 & $0{,}000000000$ & $0{,}00000000000$ & $-1{,}39\cdot10^{-17}$ & підтверджено \\
Табл. 13: Білок, г: приріст $C^*$, 1944 & $0{,}000000000$ & $0{,}00000000000$ & $0$ & підтверджено \\
Табл. 13: Кальцій, г: приріст $C^*$, 1939 & $0{,}000253902$ & $0{,}00025390171$ & $-2{,}92\cdot10^{-10}$ & підтверджено \\
Табл. 13: Кальцій, г: приріст $C^*$, 1944 & $0{,}000226499$ & $0{,}00022649904$ & $3{,}94\cdot10^{-11}$ & підтверджено \\
Табл. 13: Залізо, мг: приріст $C^*$, 1939 & $0{,}000000000$ & $0{,}00000000000$ & $1{,}39\cdot10^{-17}$ & підтверджено \\
Табл. 13: Залізо, мг: приріст $C^*$, 1944 & $0{,}000000000$ & $0{,}00000000000$ & $-2{,}78\cdot10^{-17}$ & підтверджено \\
Табл. 13: Вітамін A, тис. МО: приріст $C^*$, 1939 & $0{,}000020012$ & $0{,}00002001164$ & $-3{,}64\cdot10^{-10}$ & підтверджено \\
Табл. 13: Вітамін A, тис. МО: приріст $C^*$, 1944 & $0{,}000036224$ & $0{,}00003622423$ & $2{,}27\cdot10^{-10}$ & підтверджено \\
Табл. 13: Вітамін B1, мг: приріст $C^*$, 1939 & $0{,}000000000$ & $0{,}00000000000$ & $1{,}39\cdot10^{-17}$ & підтверджено \\
Табл. 13: Вітамін B1, мг: приріст $C^*$, 1944 & $0{,}000000000$ & $0{,}00000000000$ & $-2{,}78\cdot10^{-17}$ & підтверджено \\
Табл. 13: Вітамін B2, мг: приріст $C^*$, 1939 & $0{,}000441667$ & $0{,}00044166688$ & $-1{,}17\cdot10^{-10}$ & підтверджено \\
Табл. 13: Вітамін B2, мг: приріст $C^*$, 1944 & $0{,}000362655$ & $0{,}00036265465$ & $-3{,}48\cdot10^{-10}$ & підтверджено \\
Табл. 13: Ніацин, мг: приріст $C^*$, 1939 & $0{,}000000000$ & $0{,}00000000000$ & $-1{,}39\cdot10^{-17}$ & підтверджено \\
Табл. 13: Ніацин, мг: приріст $C^*$, 1944 & $0{,}000000000$ & $0{,}00000000000$ & $0$ & підтверджено \\
Табл. 13: Вітамін C, мг: приріст $C^*$, 1939 & $0{,}000108088$ & $0{,}00010808814$ & $1{,}37\cdot10^{-10}$ & підтверджено \\
Табл. 13: Вітамін C, мг: приріст $C^*$, 1944 & $0{,}000149432$ & $0{,}00014943201$ & $5{,}10\cdot10^{-12}$ & підтверджено \\
Підрозд. 3.1: найбільший модуль різниці приросту і 0{,}01 $b_ju_j^*$ менший за $10^{-12}$ (18 випадків) & $<1{,}00\cdot10^{-12}$ & $2{,}78\cdot10^{-17}$ & --- & підтверджено (нерівність виконується) \\
Підрозд. 3.1: $C^*$ за всіх вимог +1\%, 1939, дол./день & $0{,}109748900989$ & $0{,}10974890098882$ & $-1{,}76\cdot10^{-13}$ & підтверджено \\
Підрозд. 3.1: $C^*$ за всіх вимог +1\%, 1944, дол./день & $0{,}165751934533$ & $0{,}16575193453294$ & $-6{,}03\cdot10^{-14}$ & підтверджено \\
\end{longtable}}

\end{document}
